\section{Methodology}
\label{sec:methodology}

We propose an agentic retrieval-augmented generation framework that
addresses three failure modes endemic to single-pass retrieval pipelines:
mis-targeted retrieval strategies, generation grounded in irrelevant
context, and unverified hallucinated content. The framework is realised
as a deterministic finite-state controller orchestrating five
specialised components: an adaptive router, a multi-path retriever, a
sufficiency grader, a generator, and a hallucination verifier. Two
auxiliary components, a query rewriter and a decomposer, are invoked
under failure-driven escalation. The complete control flow is bounded
by hard iteration limits, ensuring tractable worst-case latency and
cost.

\subsection{System Overview}
\label{sec:overview}

Let $q$ denote a user query and $\mathcal{C}$ a corpus indexed under a
hybrid representation comprising dense passage embeddings, sparse
lexical signals, and a knowledge graph $\mathcal{G}=(V,E)$ derived from
entity and relation extraction over $\mathcal{C}$. The system maintains
a state vector
$s = (q, q', \pi, K, a, \gamma, U, c_r, c_v)$
where $q'$ is the active retrieval query (initially $q'=q$), $\pi$ is
the active retrieval profile, $K$ is the accumulated chunk set, $a$ is
the candidate answer, $\gamma \in \{0,1\}$ is the grounding indicator,
$U$ is the set of ungrounded claims, and $c_r, c_v$ are independent
counters bounding the retrieval and verification loops respectively.
State transitions are governed by a graph $\mathcal{T}$ in which each
node corresponds to a single component invocation and each edge encodes
a deterministic decision rule.

\subsection{Adaptive Query Routing}
\label{sec:router}

The router maps an incoming query to a retrieval profile
$\pi \in \Pi = \{\textsc{precise}, \textsc{semantic}, \textsc{local},
\textsc{multihop}\}$. Each profile encodes a distinct combination of
retrieval paths and rank-fusion weights tailored to a class of query
intents: lexical exact match, dense semantic similarity, single-entity
graph neighbourhood, and multi-entity cross-document reasoning. A fifth
profile, \textsc{full}, executes the union of all paths but is excluded
from the router's choice set; it is reserved for terminal escalation
described in Section~\ref{sec:escalation}.

Routing is performed by a language model classifier conditioned on a
prompt that enumerates each profile, its activating signals, and
representative examples. The classifier emits a JSON object
$\{\pi, \rho\}$ where $\rho \in [0,1]$ is a self-reported confidence.
When $\rho$ falls below a threshold $\tau_\rho$, the router falls back
to the \textsc{semantic} profile, which functions as the modal default
across query distributions observed in pilot deployments. To amortise
classifier cost, decisions are cached in a fixed-capacity LRU keyed by
the SHA-256 prefix of the normalised query concatenated with a hash of
the classifier prompt template. Cache entries carry a tri-state outcome
field updated by downstream signals: a successful verification marks
the entry as confirming, while terminal failure marks it as discrediting
and excludes the cached profile from subsequent classifications of the
same query.

\subsection{Multi-Path Retrieval with Weighted Reciprocal Rank Fusion}
\label{sec:retrieval}

Each profile $\pi$ specifies a set of retrieval paths
$P_\pi \subseteq \mathcal{P}$ and a weight function
$w_\pi: \mathcal{P} \to \mathbb{R}_{>0}$. Paths execute concurrently
under a semaphore that bounds simultaneous load on shared backends.
Failed paths are dropped silently provided at least one path returns
results; complete failure raises a controlled exception that triggers
the terminal insufficiency state described in Section~\ref{sec:termination}.

Per-path ranked outputs are fused via weighted reciprocal rank fusion.
For a chunk $d$ appearing at rank $r_p(d)$ in path $p$,
\begin{equation}
S(d) = \sum_{p \in P_\pi} w_\pi(p) \cdot \frac{1}{k + r_p(d)}
\end{equation}
where $k$ is a stability constant. The fused list is filtered by a
minimum score threshold $\tau_{\text{rrf}}$ to remove the long tail of
low-confidence candidates, then truncated to a profile-specific cap
$\kappa$. The truncated set is reranked by a cross-encoder model whose
pairwise scores yield a final ordering, after which a relevance
threshold $\tau_{\text{rerank}}$ removes residual noise. The resulting
chunk set $K$ satisfies $|K| \le \chi$ where $\chi$ is the global
context budget. The cap $\kappa$ is chosen such that reranker batches
fit within available accelerator memory, eliminating the need for
runtime batch-size backoff.

\subsection{Sufficiency Assessment}
\label{sec:grader}

Prior to generation, a grader assesses whether the retrieved set $K$
contains adequate evidence to answer $q$. The grader is implemented as
a single language model invocation whose prompt comprises the full text
of all chunks in $K$ followed by a task-specific suffix. Output is a
binary judgment $\sigma \in \{0, 1\}$ accompanied by a one-sentence
diagnostic rationale used by downstream rewriting. We denote the
grader's mapping by $G(q, K) \mapsto (\sigma, \delta)$ where $\delta$ is
the diagnostic.

Critically, the grader, generator, and hallucination verifier share an
identical prompt prefix consisting of the full chunk corpus. Under
inference engines supporting automatic prefix caching, this enables the
attention KV cache computed during the grader's forward pass to be
reused by the subsequent generator and verifier invocations,
amortising the prefill cost of the largest prompt segment across three
calls. A grader yielding $\sigma = 1$ transitions the controller to
generation; otherwise the system enters the iterative refinement
regime.

\subsection{Iterative Refinement and Failure-Driven Escalation}
\label{sec:escalation}

When the grader returns insufficiency, the controller selects a
remediation strategy as a deterministic function of the retrieval
counter $c_r$. We adopt a two-stage escalation policy motivated by
distinct hypotheses about the cause of failure.

At $c_r = 0$, the system invokes a query rewriter that conditions on
the original query and the grader's diagnostic to produce a reformulated
query $q'$. The rewriter targets surface-level mismatches between
user vocabulary and corpus vocabulary, including synonym substitution,
disambiguation of polysemous terms, and decomposition of compound noun
phrases. Crucially, the active profile $\pi$ is held constant: the
hypothesis under test is that the query is poorly formulated, not that
the retrieval strategy is wrong.

At $c_r = 1$, the rewriter has failed to elicit sufficient evidence,
permitting the controller to discard the query-formulation hypothesis.
The system then invokes a decomposer that partitions the original
query into a small set of independent sub-queries
$\{q^{(1)}, \ldots, q^{(m)}\}$ where $m \le 4$. Each sub-query is
retrieved concurrently under the \textsc{full} profile, which activates
the union of all retrieval paths to maximise recall. Sub-query results
are fused, deduplicated, and capped to prevent unbounded growth of the
context set. This cycle constitutes the system's terminal retrieval
attempt; subsequent failure triggers controlled abstention rather than
further escalation.

\subsection{Grounded Generation}
\label{sec:generation}

When sufficiency is established, the generator produces a candidate
answer $a$ conditioned on $K$ and $q$ via a prompt that constrains the
model to draw exclusively from the provided evidence. The prompt
includes an explicit honesty directive instructing the model to declare
inability when the evidence is inadequate, rather than fabricating.
This directive interacts productively with the downstream verifier:
declarations of inability contain no factual claims and therefore pass
verification trivially, providing a low-overhead path for genuinely
unanswerable queries.

\subsection{Hallucination Verification}
\label{sec:verification}

The verifier examines $a$ against $K$ and emits a binary grounding
judgment $\gamma \in \{0, 1\}$ together with the set
$U = \{u_1, \ldots, u_\ell\}$ of claims appearing in $a$ but unsupported
by $K$. We deliberately reject continuous grounding scores and
threshold-based partial acceptance. The motivating assumption is that
in regulated application domains, even a single fabricated claim can
propagate into downstream decisions with high cost, and post-hoc
filtering of partial fabrications cannot reliably separate the
unsupported fragment from the surrounding text.

When $\gamma = 1$, the answer is returned with high confidence. When
$\gamma = 0$ and the verification counter satisfies $c_v < c_v^{\max}$,
the system invokes a targeted retriever that uses $U$ as a query to
augment $K$ with evidence directly addressing the unsupported claims.
Generation is then re-executed against the enlarged context. This
inner loop is bounded by $c_v^{\max}$ and accumulates retrieved
chunks across iterations, subject to a deduplication and capping rule.

\subsection{Termination and Confidence Reporting}
\label{sec:termination}

The controller admits exactly two terminal failure states. The first,
denoted $\textsc{end}_{\text{insufficient}}$, is entered when either
the retrieval loop exhausts its iteration budget without achieving
sufficiency, or the verification loop exhausts its budget without
achieving grounding. In both cases the system returns no answer; the
generator output, if any, is suppressed to prevent propagation of
unverified content. The second, denoted $\textsc{end}_{\text{degraded}}$,
fires only when the verifier itself fails due to an infrastructure
fault such as a malformed model response. In this case the generator's
last output is returned with a confidence indicator marking the answer
as unverified, allowing calling systems to apply domain-specific
filtering policies.

System output is therefore a triple $(a^\star, \gamma^\star, U^\star)$
augmented with a confidence label
$\zeta \in \{\textsc{high}, \textsc{degraded}, \textsc{none}\}$ and a
trace record sufficient for post-hoc analysis. This separation between
inability and uncertainty enables downstream consumers, including human
reviewers and automated decision pipelines, to apply differentiated
handling: a $\zeta = \textsc{none}$ result is a categorical refusal and
should never be presented as an answer, whereas a $\zeta = \textsc{high}$
result carries the system's strongest grounding guarantee.

\subsection{Cost and Latency Bounds}
\label{sec:bounds}

The framework's worst-case behaviour is fully characterised by the
iteration limits $c_r^{\max}$ and $c_v^{\max}$. Let $L_R$, $L_G$,
$L_V$, $L_W$, and $L_D$ denote the language model invocation counts
for the router, grader, verifier, rewriter, and decomposer respectively,
each equal to one per invocation, and let $L_E$ denote the generator
count. The total invocation count for a query traversing the worst
admissible path is
\begin{equation}
L_{\text{worst}} = L_R + (c_r^{\max}+1) L_G + L_W + L_D + (c_v^{\max}+1)(L_E + L_V)
\end{equation}
which under the chosen limits $c_r^{\max}=2$ and $c_v^{\max}=1$ yields
nine model invocations. Empirically observed latencies on the optimal
path lie within the range required for interactive use, while pathological
queries terminate cleanly within a bounded multiple thereof. The hard
upper bound on cost, in conjunction with the explicit abstention
mechanism, distinguishes this framework from open-ended agent loops
whose cost is governed by emergent heuristics rather than analytic
guarantees.
